% Pro M. Fonteio (pro-fonteio) --- English translation
% Translated by Claude (Anthropic) from the Latin (Perseus, Clark text).
% Style guide: STYLE.md.
%
% A defence in 69 BC of M. Fonteius, former governor of Transalpine
% Gaul (74-72 BC), prosecuted for extortion by his Gallic subjects.
% The manuscript is mutilated; the speech survives only in fragments.
% The opening sections are damaged.

\ciceroSection{1} \textit{[The opening is lost.]} ...or to have paid, or so to have settled as all others have settled? For thus I defend Marcus Fonteius, gentlemen, and so contend that, after the Valerian law was passed, when you were quaestor, up to the quaestorship of Titus Crispinus, no one paid otherwise; that he followed the authority of all his predecessors, and that all who came after followed his.

\ciceroSection{2} What do you accuse, what do you reproach? For when, in the dodrans and quadrans tablets which he says were instituted by Hirtuleius, he longs for an office of Fonteius, I cannot reckon whether he himself errs or wishes to lead you into error. For I ask of you, Marcus Plaetorius: can our cause be approved by you yourself, if, in the matter in which Marcus Fonteius is accused by you, he has as authority that very Hirtuleius whom you most praise; while in that matter in which you praise Hirtuleius, Fonteius is found to have done the same? You reproach the kind of payment; the public records show Hirtuleius paid in the same way. You praise him because he set up the dodrans tablets; the same Fonteius set up, and of the same kind of money. For lest perhaps you should be ignorant and reckon that those tablets had to do with a different reckoning of old debt --- they were set up for one cause and in one kind. For with the publicans who held Africa, who held the Aquileian harbour duty\ldots \textit{[The text breaks off.]}

\ciceroSection{3} \textit{[Resuming.]} ...no one, no one, I say, gentlemen, will be found who shall say that he gave one penny to Marcus Fonteius in his quaestorship, or that the latter took from that money which was paid for the treasury. In the records of no man will any sign of this theft be found, no trace of waste or diminution among those entries. And yet men, if any have been accused on this kind of inquiry, we see have been reproached first by witnesses; for it is hard that he who has given money to a magistrate should not either be led by hatred to speak or compelled by religious obligation. Then if by some favour the witnesses are deterred, the records at least remain uncorrupted and whole. Suppose all were most friendly to Fonteius, suppose so great a number of men --- most unknown and most foreign --- spared his life, looked out for his standing. The matter itself, however, and the reckoning of letters, and the drawing-up of the records have this force: that out of receipts and outlays, whatever is feigned, or stolen, or does not stand, appears. They have all entered moneys received by the Roman people; if they at once paid out or gave to others sums equally great, so that what is received by the Roman people is the outlay to someone, surely nothing can have been taken away. But if they have brought anything home, out of their chests, out of\ldots \textit{[Text broken.]}

\ciceroSection{4} ...the faith of gods and men! No witness is found in 23 million sesterces! Of how many men? Of more than six hundred. In what lands was the business done? In that, in that place, I say, which you see. Was money given out of order? Indeed not a single penny was moved without many letters. What then is this prosecution which can more easily climb the Alps than a few steps of the treasury, which more diligently defends the treasury of the Ruteni than of the Roman people, which more willingly uses unknown than known witnesses, foreign than domestic; which thinks the charge can be more plainly confirmed by the lust of barbarians than by the letters of our men?

\ciceroSection{5} The reckoning of two magistracies, in each of which he was occupied with handling and managing the greatest sums of money, the triumvirate and the quaestorship, is so rendered, gentlemen, that in those things which were done before our eyes, which had to do with many men, which were drawn up in public and private records, no sign of theft, no suspicion of any offence is found.

\ciceroSection{6} The Spanish legateship followed at a most troubled time of the commonwealth, when, with the coming of Lucius Sulla and the greatest armies into Italy, the citizens differed by force, not by trials and laws; and in this state of the commonwealth despaired of, what kind of\ldots \textit{[Text broken.]}

\ciceroSection{7} If no money has been counted out, of what money is the fiftieth?

\ciceroSection{8} The greatest quantity of grain from Gaul; the most ample troops of foot from Gaul; horsemen most in number from Gaul.

\ciceroSection{9} That afterwards the Gauls were going to drink it more diluted --- which they were going to reckon as poison.

\ciceroSection{10} That Plaetorius's mother, while she lived, kept a school; after she died, she had schoolmasters.

\ciceroSection{11} That under this praetor Gaul was overwhelmed by debt. From whom do they say loans of such great sums were made? From the Gauls? Nothing less. From whom then? From the Roman citizens who do business in Gaul. Why do we not hear their words? Why are no records of theirs brought forth? I press the prosecutor still and stand by him, gentlemen; I press him, I say, and demand the witnesses. I spend more labour and toil in this case in demanding witnesses than other defenders in refuting them. Boldly I say this, gentlemen, I do not rashly affirm it: Gaul is stuffed with businessmen, full of Roman citizens. No Gaul does any business without a Roman citizen. No coin in Gaul is moved without the records of Roman citizens.

\ciceroSection{12} See to what I come down, gentlemen, how far I seem to depart from my custom and from caution and diligence. Let one set of records be brought forth in which there is some trace which signifies that money was given to Marcus Fonteius. Let them produce one out of so great a number of businessmen, settlers, publicans, farmers, cattle-men, as witness; I shall grant the prosecution to be true. By the immortal gods! What is this case, what defence? Marcus Fonteius governed the province of Gaul, which is made up of those kinds of men and states which (to leave aside the old) partly within our own memory waged bitter and long-lasting wars with the Roman people; partly were lately subdued by our generals, lately overcome by war, lately marked by triumphs and monuments, lately fined of fields and cities by the senate; partly with Marcus Fonteius himself contended with iron and hand, and by his much sweat and toil fell under the imperium and rule of the Roman people.

\ciceroSection{13} In the same province is Narbo Martius, a colony of our citizens, the watchtower and bulwark of the Roman people set against and opposed to those very nations. There is also the city of Massilia, of which I have spoken before --- of bravest and most faithful allies, who have made up the perils of the Gallic wars by special rewards from the Roman people. There are besides the largest number of Roman citizens and knights, most honourable men. Over this province, which from this variety of kinds was made up, Marcus Fonteius (as I said) governed. Those who were enemies he subdued; those who had lately been so, he compelled to depart from those fields of which they had been fined. On the rest, who had often been overcome in great wars on this account that they should always obey the Roman people, he commanded great cavalries for those wars which were then being waged in the whole world by the Roman people; great sums of money for their pay; the greatest quantity of grain for sustaining the Spanish war.

\ciceroSection{14} He who carried this through is called into court; you who were not present at the matter learn the case together with the Roman people. Those speak against him to whom most unwillingly it was commanded; those speak who were compelled to depart from their fields by Gnaeus Pompeius's decree; those speak who out of the slaughter and flight of war now first dare to stand up against an unarmed Marcus Fonteius. What? The Narbonensian colonists --- what do they wish, what do they think? They wish him kept safe through you; they think themselves kept safe through him. What of the state of the Massilians? Him present they have furnished with those most ample honours which they had; you absent they ask and beseech that their religious obligation, their praise, their authority may seem to have had some weight before your minds.

\ciceroSection{15} What? What is the will of the Roman citizens? There is no one out of so great a number who does not reckon him to have deserved most excellently of the province, of the imperium, of the allies and of the citizens. Since therefore, gentlemen, you have learned who attacked Marcus Fonteius, who wish him defended, decide now what your equity, what the dignity of the Roman people demands: whether you would prefer to trust and to consult for your colonists, your businessmen, your most friendly and most ancient allies, than for those to whom you ought neither (on account of their anger) to give faith, nor (on account of their faithlessness) to give honour.

\ciceroSection{16} What? If I should bring even a greater abundance of most honourable men which can be a witness to his virtue and innocence, will the agreement of the Gauls yet have more weight than the highest authority of all? When Fonteius was governing Gaul, you know, gentlemen, that the largest armies of the Roman people were in the two Spains, and the most distinguished generals. How many Roman knights, how many tribunes of the soldiers, of what kind and how many and how often legates went out to them! Besides, the largest and most adorned army of Gnaeus Pompeius wintered in Gaul under Marcus Fonteius's command. Does fortune itself seem to you to have wished there to be witnesses and partners enough, fit enough, of those things which were done in Gaul under Marcus Fonteius as praetor? Whom can you produce out of so great a number of men in this case as witness? Who is there of that number whom you would approve as authority? That man we shall now use as praiser and witness.

